\item \textbf{Problema de repaso.} A medida que una onda sonora pasa a través de un gas, las compresiones son tan rápidas o tan separadas que la conducción térmica se evita por un intervalo de tiempo despreciable o por un grosor de aislamiento efectivo. Las compresiones y enrarecimientos son adiabáticos. 
(a) Demuestre que la rapidez del sonido en un gas ideal es
$$ v = \sqrt{\frac{\gamma RT}{M}} $$
donde $M$ es la masa molar. La rapidez del sonido en un gas está dada por la ecuación 4.35; use esa ecuación y la definición del módulo volumétrico. 
(b) Calcule la rapidez teórica del sonido en el aire a $20^\circ\text{C}$ y establezca cómo se compara con el valor en la tabla 8.1. Tome $M = 28.9\text{ g/mol}$. 
(c) Demuestre que la rapidez del sonido en un gas ideal es
$$ v = \sqrt{\frac{\gamma k_B T}{m_0}} $$
donde $m_0$ es la masa de una molécula. 
(d) Establezca cómo el resultado en el inciso (c) se compara con las magnitudes de velocidad más probable, promedio y rms moleculares.
